\section{Completeness of the bounded census}\label{app:ordinary-census}
This appendix supplies the mathematical coverage argument for the finite
search. Checking that emitted tensors are valid and pairwise nonisomorphic
does not by itself prove that every ring occurs. We therefore give the
complete presentations, the elimination branches, and the exact acceptance
condition. All parameters range over integers in $[0,M]$, where $M=31$;
write $\Lambda=1000$ for the exclusive dimension bound.

\subsection{Complete presentations}
By Lemma~\ref{core:lem:rank-four-algebra}, every rank-four fusion ring is
commutative. Duality on its three nonunit basis elements is either the
identity or a transposition. Reciprocity gives the following two
presentations, with $N_0=I$ and our row convention
$(N_u)_{vw}=N_{uv}^{w}$.

In the self-dual case use the matrices and six associativity equations
of Section~\ref{finite:sec:census}, with nonnegative parameters
$(a,b,c,d,e,f,g,h,i,j)$.

In the case with duality $(0)(1)(23)$ the parameters are $(a,b,c,d,e,f)$,
and
\[
 N_1=\begin{pmatrix}0&1&0&0\\1&a&b&b\\0&b&c&d\\0&b&d&c\end{pmatrix},\quad
 N_2=\begin{pmatrix}0&0&1&0\\0&b&c&d\\0&d&e&f\\1&c&e&e\end{pmatrix},\quad
 N_3=N_2^{\mathsf T}.
\]
Their commutativity is equivalent to
\begin{align}
 ac-b^2+2be-c^2-d^2+1&=0,\notag\\
 ad-b^2+be+bf-2cd&=0,\label{census:eq:ns-assoc}\\
 b(c-d)+d(e-f)&=0,\notag\\
 -c^2+d^2-e^2+f^2-1&=0.\notag
\end{align}
In both cases commutativity of the displayed matrices suffices for
associativity of the proposed multiplication. Indeed, set
$E=N_uN_v-\sum_w N_{uv}^{w}N_w$. Its zeroth row vanishes, and it
commutes with every $N_t$. Its row $t$ is the zeroth row of
$N_tE=EN_t$, hence also vanishes. Thus $E=0$, proving the full
regular multiplication identities. The unit, duality, integrality,
nonnegativity, and reciprocity axioms are already built into the
presentations.

\subsection{The self-dual branches}
Permutations of the three nonunit elements act transitively on the six
mixed coordinates $b,c,d,f,h,i$. Every isomorphism class therefore
has a representative with $b=\min(b,c,d,f,h,i)$. We enumerate all
such representatives before imposing the final isomorphism reduction.

\paragraph{Positive minimal mixed coefficient.}
Suppose $b>0$. Enumerate $b,c,d,e,f$, with $c,d,f\geq b$,
and use $H_0,G_0,J_0,A,C,B_0$ from Section~\ref{finite:sec:census}.
The first, second, and fourth associativity equations give the displayed
formulas there for $h,g,j$; the third gives
$Ca+Ai=-B_0$, namely \eqref{finite:eq:census-linear}.
If $A\ne0$, put $q=\gcd(A,C)$. There are no solutions unless
$q\mid B_0$; otherwise $a$ lies in the unique residue class solving
\[
 (C/q)a\equiv-B_0/q\pmod{|A/q|},
\]
and $i$ is determined by \eqref{finite:eq:census-linear}. A modulus of one
imposes no restriction. If $A=0\ne C$, determine $a$ and
retain every bounded $i$. If $A=C=0$, require $B_0=0$ and
retain every bounded pair $(a,i)$. In all cases retain only integral
$g,h,j$, with $a,g,j\in[0,M]$ and $h,i\in[b,M]$, and test all
six original equations.

The implementation discards outer tuples only by necessary interval
conditions. Besides $b,c,d,f\in[b,M]$ and $e\in[0,M]$, these are
\begin{gather*}
 b(b+c)\leq H_0\leq M(e+b+c),\qquad
 bc\leq G_0\leq M(d+b+c),\\
 b^2\leq J_0\leq M(f+b+c),\\
 \min(0,CM)+\min(Ab,AM)\leq-B_0
 \leq\max(0,CM)+\max(Ab,AM).
\end{gather*}
The first three follow from the displayed equations for $h,g,j$;
the last bounds the left side of \eqref{finite:eq:census-linear} over its
allowed rectangle. The further necessary dimension inequality
\[
 2+2b^2+2c^2+d^2+2e^2+f^2<\Lambda
\]
is the bound on $B_{11}$ after discarding its nonnegative summand $a^2$.
Here $B=\sum_u N_uN_u^{\mathsf T}$ is the Casimir matrix.

\paragraph{One zero mixed coefficient.}
Suppose $b=0<c$. The equations imply
\[
 h=\frac{d^2+e^2-ad-1}{c},\qquad
 i=\frac{e(d+f-a)}{c},\qquad
 j=\frac{c^2+e^2+f^2-af-1}{c},
\]
and $eg=h(d-f)+ei-cd$.
For $e>0$, enumerate $c,d,e,f$ and the progression
$a\equiv d+f\pmod{c/\gcd(c,e)}$. The bounds on $h,i,j$ give
necessary bounds on $a$; exact division determines $g$. Test all
six equations and all parameter bounds.

For $e=0$, one has $i=0$, and the first and third equations imply
$d,h>0$ and $\gcd(c,d)=1$. Indeed $d=0$ would give $ch=-1$,
and $h(d-f)=cd>0$ forces $h>0$; divisibility of $d(d-a)-1$
by $c$ gives coprimality. Enumerate $c,d$ and
$a\equiv d-d^{-1}\pmod c$, with the congruence unrestricted for
$c=1$. Determine $h$, put $f=d-cd/h$, and determine $j$ by the
displayed formula. Retain every bounded $g$, testing exact integrality,
all bounds, and all six equations.

\paragraph{Two zero mixed coefficients, with $e>0$.}
For $b=c=0<e$, the first, second, and fourth equations yield
$a=d+f$ and $df=e^2-1$. Enumerate every bounded factor pair,
including the zero-factor possibilities when $e=1$. For every
bounded $h,i$, determine $g,j$ from
\[
 eg=h(d-f)+ei,\qquad ej=eh+i(f-d).
\]
Retain the integral bounded solutions satisfying all six equations.

\paragraph{The remaining zero branch.}
For $b=c=e=0$, the equations force $a=0$ and $d=f=1$ and leave
\[
 gi+hj=h^2+i^2-2.
\]
For $h>0$, enumerate $h,i,g$ and determine $j$; for $h=0<i$,
determine $g$ and retain every bounded $j$. The case $h=i=0$
is impossible. As before, test all original equations and bounds.

These four branches are exhaustive. In each division step a nonzero
denominator has been explicitly assumed; when a coefficient vanishes,
every remaining bounded solution is retained. Finally keep the
lexicographically least tuple among its permutations preserving the
minimal value of $b$. Each orbit has exactly one such tuple. Thus the
self-dual enumeration has neither an omitted zero case nor a duplicate
based-ring isomorphism class.

\subsection{The branch with a dual pair}
Enumerate $c,d,e$. The last equation in \eqref{census:eq:ns-assoc} determines
the unique possible nonnegative integer
$f=\sqrt{1+c^2-d^2+e^2}$; test its square and bound exactly.
If $c\ne d$, the third equation determines
$b=-d(e-f)/(c-d)$. If $c=d$, require $d(e-f)=0$ and retain
every bounded $b$. Determine $a$ from the first equation when $c>0$,
from the second when $c=0<d$, and retain every bounded $a$ when
$c=d=0$. Test all four original equations and bounds.
The only unit-preserving relabelling compatible with this duality
exchanges the dual pair. It leaves the six parameters unchanged, so
these tuples already represent distinct isomorphism classes.

\subsection{The dimension test and the coverage conclusion}
The positive dimension vector is an eigenvector of the symmetric
nonnegative Casimir matrix $B$, with eigenvalue $D=\FPdim(R)$.
It follows that $D$ is its largest eigenvalue. Therefore
\begin{equation}\label{census:eq:pd}
 D<\Lambda\quad\Longleftrightarrow\quad
 \Lambda I-B\text{ is positive definite}
 \quad\Longleftrightarrow\quad
 \det(\Lambda I-B)_{[1,k]}>0\ (1\leq k\leq4).
\end{equation}
Here $[1,k]$ denotes the leading principal $k$-by-$k$ submatrix.
The final equivalence is Sylvester's criterion. All determinants are
integers; a zero determinant is rejected, so equality at the boundary
is never admitted. The C++ enumeration uses fraction-free elimination,
and the Python tensor checker recomputes the same quantities with
arbitrary-precision integers.

\begin{proposition}\label{census:prop:coverage}
The preceding procedure with $M=31$ and $\Lambda=1000$ emits exactly
one representative of every rank-four fusion ring with $D<1000$.
\end{proposition}
\begin{proof}
Every rank-four fusion ring has one of the two complete presentations.
All solutions of their associativity equations in the parameter box
occur in the enumerated branches, and the isomorphism reduction retains
exactly one representative. The acceptance condition
\eqref{census:eq:pd} is equivalent to the desired strict dimension
bound. Conversely, every ring below that bound has $\mu<\sqrt D$
by Lemma~\ref{core:lem:coefficient-bound}, hence $\mu\leq31$, and therefore
lies in the searched box. This proves both directions.
\end{proof}

The executable output has $\CensusSelfDual$ self-dual classes and
$\CensusNonSelfDual$ classes with a dual pair, for a total of $\CensusTotal$.
Completeness follows from the proposition, not merely from these totals.
As an additional implementation check, removing both dimension-pruning
steps gives $16065$ classes at multiplicity at most $31$; an independently
implemented strict dimension test retains precisely the same $\CensusTotal$
parameter tuples. A full small-box test examines all
$3^{10}+3^6=59778$ assignments at multiplicity at most two and recovers
exactly the same $27$ classes. The unpruned experiment shares the
Diophantine elimination; the small-box test avoids it but shares the
tensor presentations and canonicalization. Neither is asserted to be a
second independent proof of the complete enumeration algorithm.
